\setcounter{section}{14}

  \zwischenueberschrift{Soal latihan}

\inputaufgabe
{}
{

Misalkan $M$ sebuah
\definitionsverweis {manifold diferensiabel}{}{}
dengan
\definitionsverweis {bundel kotangen}{}{}

\mathl{T^*M}{.} Tunjukkan bahwa pada
\mathl{\bigwedge^k T^*M}{} untuk setiap $k$ dapat didefinisikan suatu
\definitionsverweis {topologi}{}{}
sedemikian sehingga, untuk setiap peta
\maabb {\alpha} { U } { V
} {}
pemetaan
\maabbdisp {} {\bigwedge^k T^*U} { \bigwedge^k T^*V \cong V \times \bigwedge^k \R^{n*}
} {}
merupakan sebuah
\definitionsverweis {homeomorfisme}{}{.}

}
{} {}
Dengan demikian, kita dapat berbicara tentang bentuk diferensial kontinu maupun terukur.

\inputaufgabe
{}
{

Misalkan $M$ sebuah
$C^2$-\definitionsverweis {manifold diferensiabel}{}{}
dengan
\definitionsverweis {bundel kotangen}{}{}

\mathl{T^*M}{.} Tunjukkan bahwa
\mathl{\bigwedge^k T^*M}{} merupakan sebuah manifold diferensiabel untuk setiap $k$.

}
{} {}
Dengan demikian, kita juga dapat berbicara tentang bentuk diferensial yang terdiferensialkan secara kontinu. Secara lebih umum, jika $M$ sebuah
$C^m$-\definitionsverweis {manifold}{}{}
dan

\mavergleichskette
{\vergleichskette
{  \ell
}
{ < }{ m
}
{  }{
}
{    }{
}
{   }{
}
}
{}{}{}
maka
\mathl{{ \mathcal E }^{ k }_{\ell } ( M  )}{} menyatakan himpunan semua bentuk diferensial yang terdiferensialkan secara kontinu hingga orde $\ell$ dan berderajat $k$.

\inputaufgabe
{}
{

Misalkan $M$ sebuah
\definitionsverweis {manifold diferensiabel}{}{}
dan
\mathl{{ \mathcal E }^{ k } (  M   )}{} himpunan semua
$k$-\definitionsverweis {bentuk}{}{}
pada $M$. Tunjukkan bahwa ${ \mathcal E }^{ k } (  M   )$ merupakan sebuah
$R$-\definitionsverweis {modul}{}{}
untuk

\mavergleichskette
{\vergleichskette
{R
}
{ = }{ C^0(M,\R)
}
{  }{
}
{    }{
}
{   }{
}
}
{}{}{.}

}
{} {}

\inputaufgabe
{}
{

Misalkan $M$ sebuah
\definitionsverweis {manifold diferensiabel}{}{,}

\mavergleichskette
{\vergleichskette
{ P
}
{ \in }{ M
}
{  }{
}
{    }{
}
{   }{
}
}
{}{}{}
sebuah titik, dan

\mavergleichskette
{\vergleichskette
{ f
}
{ \in }{ C^1(M,\R)
}
{  }{
}
{    }{
}
{   }{
}
}
{}{}{}
sebuah
\definitionsverweis {fungsi yang terdiferensialkan secara kontinu}{}{.}
Misalkan

\mavergleichskette
{\vergleichskette
{ v
}
{ \in }{ T_PM
}
{  }{
}
{    }{
}
{   }{
}
}
{}{}{}
sebuah
\definitionsverweis {vektor singgung}{}{,}
yang direpresentasikan oleh suatu lintasan diferensiabel
\maabbdisp {\gamma} {{]{-\delta},\delta [}} {M
} {}
dengan

\mavergleichskette
{\vergleichskette
{ \gamma(0)
}
{ = }{ P
}
{  }{
}
{    }{
}
{   }{
}
}
{}{}{.}
Tunjukkan kesamaan

\mavergleichskettedisp
{\vergleichskette
{(df)(P,v)
}
{ =} { (f \circ \gamma)'(0)
}
{ } {
}
{ } {
}
{ } {
}
}
{}{}{.}

}
{} {}

\inputaufgabegibtloesung
{}
{

Misalkan $M$ sebuah
\definitionsverweis {manifold diferensiabel}{}{}
dan
\maabbdisp {f} { M } {\R
} {}
sebuah
\definitionsverweis {fungsi yang terdiferensialkan secara kontinu}{}{.}
Pada titik

\mavergleichskette
{\vergleichskette
{ P
}
{ \in }{ M
}
{  }{
}
{    }{
}
{   }{
}
}
{}{}{}
fungsi tersebut mempunyai suatu
\definitionsverweis {ekstremum lokal}{}{.}
Tunjukkan bahwa

\mavergleichskettedisp
{\vergleichskette
{ (df)_P
}
{ =} { 0
}
{ } {
}
{ } {
}
{ } {
}
}
{}{}{.}

}
{} {}

\inputaufgabegibtloesung
{}
{

Misalkan $M$ sebuah
\definitionsverweis {manifold diferensiabel}{}{}
\definitionsverweis {kompak}{}{}
dan
\maabbdisp {f} { M } {\R
} {}
sebuah
\definitionsverweis {fungsi yang terdiferensialkan secara kontinu}{}{.}
Tunjukkan bahwa terdapat sedikitnya dua titik pada $M$
\zusatzklammer {dengan syarat $M$ mempunyai sedikitnya dua titik} {} {}
di mana
$1$-\definitionsverweis {bentuk diferensial}{}{}
$df$ bernilai nol.

}
{} {}

\inputaufgabe
{}
{

Misalkan
\mathl{i: M\subseteq \R^n}{} sebuah
\definitionsverweis {submanifold tertutup}{}{.}
Tunjukkan bahwa untuk setiap
\definitionsverweis {fungsi diferensiabel}{}{}
\maabbdisp {f} {\R^n } {\R
} {}
berlaku hubungan

\mavergleichskettedisp
{\vergleichskette
{ i^* (df)
}
{ =} { d(f \circ i )
}
{ } {
}
{ } {
}
{ } {
}
}
{}{}{.}

}
{} {}

\inputaufgabe
{}
{

Misalkan
\maabbeledisp {f} {\R} {\R
} { t } {f(t)
} {,}
sebuah
\definitionsverweis {fungsi}{}{}
\definitionsverweis {yang terdiferensialkan secara kontinu}{}{,}
dan misalkan

\mavergleichskette
{\vergleichskette
{ \omega
}
{ = }{ g(s)ds
}
{  }{
}
{    }{
}
{   }{
}
}
{}{}{}
sebuah
$1$-\definitionsverweis {bentuk diferensial}{}{}
pada $\R$. Tentukan
\mathl{f^* \omega}{.}

}
{} {}

\inputaufgabegibtloesung
{}
{

Kita meninjau pemetaan
\maabbeledisp {\varphi} {\R^3} {\R^2
} {(x,y,z)} {(xyz^2,xy-z^3)
} {.}
Hitunglah matriks pemetaan
\maabbdisp {\bigwedge^2 T_P (\varphi)} { \bigwedge^2 T_P \R^3 } { \bigwedge^2 T_{\varphi(P)}  \R^2
} {}
di titik

\mavergleichskette
{\vergleichskette
{ P
}
{ = }{ (1,3,5)
}
{  }{
}
{    }{
}
{   }{
}
}
{}{}{}
terhadap suatu basis yang sesuai.

}
{} {}

\inputaufgabe
{}
{

Kita meninjau
\definitionsverweis {pemetaan terdiferensialkan kontinu}{}{}
\maabbeledisp {\varphi} {\R^2 \setminus \{(0,0)\}} {\R^2 \setminus \{(0,0)\}
} {(u,v)} {(u^2,v^3-u)
} {,}
dan
$2$-\definitionsverweis {bentuk diferensial}{}{}

\mavergleichskettedisp
{\vergleichskette
{ \omega
}
{ =} { { \frac{   1  }{  x^2+y^2 } } dx \wedge dy
}
{ } {
}
{ } {
}
{ } {
}
}
{}{}{.}
Tentukan
\definitionsverweis {bentuk diferensial tarik balik}{}{}

\mathl{\varphi^* \omega}{.}

}
{} {}

\inputaufgabegibtloesung
{}
{

Kita meninjau
$1$-\definitionsverweis {bentuk diferensial}{}{}

\mavergleichskettedisp
{\vergleichskette
{ \omega
}
{ =} { -vdu+udv
}
{ } {
}
{ } {
}
{ } {
}
}
{}{}{}
pada
$1$-\definitionsverweis {sfera}{}{}

\mavergleichskette
{\vergleichskette
{ S^1
}
{ \subset }{ \R^2
}
{  }{
}
{    }{
}
{   }{
}
}
{}{}{,}
dengan koordinat-koordinat pada ruang ambien $\R^2$ dinotasikan sebagai
\mathkor {} {u} {dan} {v} {.}
Tentukan
\mathl{\pi^* \omega}{} di bawah
\definitionsverweis {pemetaan terdiferensialkan kontinu}{}{}
\maabbeledisp {\pi} {\R^2 \setminus \{0\}} { S^{1}
} {(x , y) } { { \frac{   1  }{  \sqrt{x^2 + y^2}  } } (x , y)
} {.}

}
{} {}

\inputaufgabegibtloesung
{}
{

Hitunglah bentuk diferensial tarik balik
\mathl{\varphi^* \tau}{} untuk

\mavergleichskettedisp
{\vergleichskette
{ \tau
}
{ =} { dx \wedge dy \wedge dz - w dx \wedge dy \wedge dw + \cos (xy)  dx \wedge dz \wedge dw-yw dy \wedge dz \wedge dw
}
{ } {
}
{ } {
}
{ } {
}
}
{}{}{}
di bawah pemetaan
\maabbeledisp {\varphi} {\R^3} {\R^4
} {(r,s,t)} { (r^2s,t, \sin  r ,e^{st}) = (x,y,z,w)
} {.}

}
{} {}

\inputaufgabegibtloesung
{}
{

Misalkan
\maabbdisp {\psi} { L } { M
} {}
sebuah
\definitionsverweis {pemetaan diferensiabel}{}{}
antara
\definitionsverweis {dua manifold diferensiabel}{}{,}
misalkan
\maabbdisp {f} { M } {\R
} {}
sebuah fungsi pada $M$, dan $\omega$ sebuah
$k$-\definitionsverweis {bentuk}{}{}
pada $M$. Tunjukkan bahwa

\mavergleichskettedisp
{\vergleichskette
{ \psi^*(f \omega)
}
{ =} { \psi^*(f)  \psi^* \omega
}
{ } {
}
{ } {
}
{ } {
}
}
{}{}{.}

}
{} {}

\inputaufgabegibtloesung
{}
{

Misalkan

\mavergleichskette
{\vergleichskette
{ W
}
{ \subseteq }{ \R^m
}
{  }{
}
{    }{
}
{   }{
}
}
{}{}{}
dan

\mavergleichskette
{\vergleichskette
{ U
}
{ \subseteq }{ \R^n
}
{  }{
}
{    }{
}
{   }{
}
}
{}{}{}
merupakan
\definitionsverweis {himpunan-himpunan bagian terbuka}{}{}
dan misalkan
\maabbdisp {\psi} { W } { U
} {}
sebuah
\definitionsverweis {pemetaan}{}{} yang
\definitionsverweis {terdiferensialkan secara kontinu}{}{.}
Misalkan
\maabbdisp {f} { U } {\R
} {}
sebuah fungsi yang terdiferensialkan secara kontinu. Simpulkan dari
aturan rantai
bahwa

\mavergleichskettedisp
{\vergleichskette
{ d (\psi^*f)
}
{ =} { \psi^*(df)
}
{ } {
}
{ } {
}
{ } {
}
}
{}{}{,}
dengan $\psi^*$ menyatakan
\definitionsverweis {penarikan balik bentuk diferensial}{}{.}

}
{} {}

  \zwischenueberschrift{Soal untuk dikumpulkan}

\inputaufgabe
{6}
{

Misalkan $M$ sebuah
\definitionsverweis {manifold diferensiabel}{}{}
dengan
\definitionsverweis {bundel kotangen}{}{}

\mathl{T^*M}{.} Misalkan $\omega$ sebuah
$k$-\definitionsverweis {bentuk diferensial}{}{,} yaitu suatu pemetaan
\maabbdisp {\omega} { M } {\bigwedge^k T^*M
} {}
dengan
\mathl{\omega(P) \in \bigwedge^k T^*_P M}{} untuk setiap
\mathl{P \in M}{,} dengan produk eksterior ini dilengkapi topologi alaminya
\zusatzklammer {lihat
Soal 14.1} {} {.}
Tunjukkan bahwa pernyataan-pernyataan berikut ekuivalen:
\aufzaehlungdrei{ $\omega$
\definitionsverweis {
kontinu}{}{.}
}{Untuk setiap
\definitionsverweis {peta}{}{}
\maabb {\alpha} { U } { V
} {}
dengan
\mathl{V \subseteq \R^n}{} dan representasi lokal
\mathl{\alpha_* \omega =  \sum_{J,\, { \# \left(   J  \right) } =k} f_J dx_J}{,}
fungsi-fungsi $f_J$ kontinu.
}{Terdapat suatu tutupan terbuka
\mathl{M= \bigcup_{i \in I} U_i}{} oleh
\definitionsverweis {domain-domain peta}{}{} $U_i$ sedemikian sehingga, dalam representasi lokal
\mathl{\alpha_{i*} \omega =  \sum_{J,\, { \# \left(   J  \right) } =k} f_{i J} dx_J}{,}
fungsi-fungsi $f_{i J}$ kontinu.
}

}
{} {}

\inputaufgabe
{4}
{

Misalkan
\maabbdisp {\varphi} { L } { M
} {}
sebuah
\definitionsverweis {pemetaan diferensiabel}{}{}
antara dua
\definitionsverweis {manifold diferensiabel}{}{}
\mathkor {} {L} {dan} {M} {.}
Misalkan
\mathkor {} {\omega \in
{ \mathcal E }^{ k } (  M   )} {dan} {\tau \in
{ \mathcal E }^{  \ell } (  M   )} {}
\definitionsverweis {bentuk-bentuk diferensial}{}{} pada $M$. Tunjukkan kesamaan

\mavergleichskettedisp
{\vergleichskette
{ \varphi^* ( \omega \wedge \tau)
}
{ =} {\varphi^* \omega \wedge \varphi^* \tau
}
{ } {
}
{ } {
}
{ } {
}
}
{}{}{.}

}
{} {}

\inputaufgabe
{4}
{

Kita meninjau
\definitionsverweis {pemetaan terdiferensialkan kontinu}{}{}
\maabbeledisp {\varphi} {\R^3 \setminus \{(0,0,0)\}} {N = \left\{ (x,y,z) \in \R^3  \mid   z \neq 0        \right\}
} {(u,v,w)} {(uvw,u^2-vw^5,u^2+v^2+w^2)
} {,}
dan
$2$-\definitionsverweis {bentuk diferensial}{}{}

\mathdisp {\omega =  z^2dx \wedge dy+{ \frac{   xy }{ z } }dx \wedge dz+(xe^y-z)dy \wedge dz} {  }

pada $N$. Tentukan
\definitionsverweis {bentuk diferensial tarik balik}{}{}

\mathl{\varphi^* \omega}{.}

}
{} {}

\inputaufgabe
{6}
{

Kita meninjau
\definitionsverweis {permukaan bola satuan}{}{}

\mavergleichskette
{\vergleichskette
{ S^2
}
{ \subseteq }{ \R^3
}
{  }{
}
{    }{
}
{   }{
}
}
{}{}{,}
dengan koordinat-koordinat pada $\R^3$ dinotasikan sebagai $x,y,z$. Untuk titik-titik manakah

\mavergleichskette
{\vergleichskette
{ P
}
{ \in }{ S^2
}
{  }{
}
{    }{
}
{   }{
}
}
{}{}{}
pembatasan
\mathkor {} {dx} {dan} {dy} {}
pada $S^2$ membentuk suatu
\definitionsverweis {basis}{}{}
dari
\definitionsverweis {ruang kotangen}{}{}

\mathl{T^*_PS^2}{.}

}
{} {}
